% TOP_PROBLEM: {"id": 3358, "title": "Lucas Numbers Modulo m", "category_id": 1, "category": {"id": 1, "name": "number_theory", "display_name": "Number Theory", "description": "Properties of integers, prime numbers, Diophantine equations.", "slug": "number-theory", "order_index": 1, "created_at": "2026-07-31T15:26:25.670Z"}, "status": "open"}
% FIRST_POSTED: null
% ENTRY_KIND: statement_only
% Imported from: batch03.tex; UnsolvedMath ID 3358
\documentclass[11pt]{article}
\usepackage[margin=25mm]{geometry}
\usepackage{fontspec}
\setmainfont{Latin Modern Roman}
\usepackage{amsmath,amssymb,mathtools}
\usepackage{unicode-math}
\setmathfont{Latin Modern Math}
\usepackage{xurl,hyperref}
\hypersetup{colorlinks=true,urlcolor=blue}
\setcounter{secnumdepth}{0}
\setlength{\parindent}{0pt}
\setlength{\parskip}{5pt}
\setlength{\emergencystretch}{3em}
\newcommand{\sref}[2]{\hyperlink{source:#1:#2}{[S#2]}}
\newcommand{\cataloguescope}[1]{\paragraph{Recorded target.}#1}
\begin{document}
% UnsolvedMath ID 3358; source catalogue code OPG-37402
\setcounter{section}{3357}
\section{Lucas Numbers Modulo m}\label{3358}
{\small\sffamily UnsolvedMath ID 3358 · Number Theory}
\subsection{Definitions and mathematical statement}

\paragraph{Shared notation.}$\mathbb N=\{1,2,\ldots\}$, and $\mathbb Z,\mathbb Q,\mathbb R,\mathbb C$ denote the integers, rationals, reals and complex numbers. Any other conventions are specified locally.


Define the Lucas sequence by $L_0=2$, $L_1=1$ and $L_n=L_{n-1}+L_{n-2}$ for $n\geq2$. For a modulus $m\geq2$, residue completeness means that for every $a\in\mathbb Z/m\mathbb Z$ there is an $n\geq0$ with $L_n\equiv a\pmod m$. The trivial modulus $1$ is residue complete automatically.
\paragraph{Statement recorded in the supplied source.}
Is the Lucas sequence residue complete modulo $m\geq2$ if and only if\[ m\in\{2,4,6,7,14\}\ \cup\ \{3^k:k\geq1\}? \] \sref{3358}{2} \sref{3358}{3}
\subsection{Short English statement}
Do Lucas numbers cover all residue classes precisely for moduli two, four, six, seven, fourteen, and positive powers of three?
\subsection{Sources}
\begin{description}
\item[\hypertarget{source:3358:1}{S1}] ULAM.AI and the UnsolvedMath Contributors, UnsolvedMath: A Curated Collection of Open Mathematics Problems, record 3358 (OPG-37402). CC BY 4.0. \url{https://huggingface.co/datasets/ulamai/UnsolvedMath}.
\item[\hypertarget{source:3358:2}{S2}] Open Problem Garden, Lucas Numbers Modulo m, node 37402, original problem page. Original question and full discussion. \url{https://www.openproblemgarden.org/op/lucas_numbers_modulo_m}.
\item[\hypertarget{source:3358:3}{S3}] Cheng Lien Lang and Mong Lung Lang, Fibonacci system and residue completeness (2013), abstract. \url{https://arxiv.org/abs/1304.2892}.
\end{description}
\label{end:3358}
\end{document}
